\subsection{An ordinary pair matched at class two in every coefficient}

\begin{result}{E.171}{Proved}{Finite reference classes determine the integral support matches}
Let \((A,B)\) be an ordinary completed pair with support \(\Sigma\).
For \(\ell\ge5\), suppose a finite class-two section \(s_\ell\)
modulo \(\ell^{N_\ell}\) has the same global abelianization,
is unramified outside \(\{\ell,\infty\}\cup\Sigma_\ell\)
for a finite subset \(\Sigma_\ell\subset\Sigma\), and matches
the actual local ordinary points at those support primes other
than \(\ell\). Suppose every prime in
\(\Sigma\setminus\Sigma_\ell\) is different from \(\ell\)
and is not congruent to \(\pm1\pmod\ell\). Put
\[
 e_\ell=\max_{q\in\Sigma_\ell,\ q\ne\ell}v_\ell(q^2-1),
 \qquad
 u_\ell=\max_{q\in\Sigma_\ell,\ q\ne\ell}v_\ell(q-1).
 \tag{E.171.1}
\]
Empty maxima are zero. If \(N_\ell\ge e_\ell+u_\ell\),
the unique global integral class-two class at coefficient
\(\ell\) matches the ordinary local point at every
\(q\in\Sigma\setminus\{\ell\}\).

For the small coefficients, suppose
\[
 \{2,23,41\}\subset\Sigma\subset
 \{2,23,41\}\cup
 \{q:q\equiv5\pmod8, q\equiv2\pmod9\}.
 \tag{E.171.2}
\]
Suppose the pair has the global abelian reduction of
\((23/64,41/64)\) modulo \(2^8\) and modulo nine, and the
actual point section of this rational base, in the corresponding
finite class-two quotients, matches its actual local ordinary
point at every support prime away from the coefficient prime.
Then one integral global two-primary class-two choice matches
all its off-diagonal support places, and one integral global
three-primary choice does so as well. These are single global
choices, not independently chosen local twists.

Consequently, if the finite-reference hypotheses hold at every
coefficient \(\ell\ge5\) which divides some \(q^2-1\),
there is one integral all-prime class-two section matching every
ordinary support point in every coefficient \(\ell\ne q\).
It has a global integral class-three lift. No matching at the
next layer or full pronilpotent lift is asserted.
\end{result}
\begin{proof}
Existence of an integral class-two lift of an ordinary pair and
the global torsor \(H^1(G_{\mathbb Q},\widehat{\mathbb Z}(2))
\simeq\mathbb Z/24\) are A.25. Fix one such lift. Choose
compatible Kummer roots so that its finite abelian cocycles are
literally those of the reference section. Their central difference
is then an ordinary continuous cocycle
\[
             z\in Z^1(G_{\mathbb Q},\mathbb Z/\ell^N(2)).
 \tag{E.171.3}
\]
For \(q\ne\ell\), its tame inertia value satisfies
\((q-1)z(T_q)=0\). At an old support prime, reduction modulo
\(\ell^e\) therefore kills inertia whenever
\(N\ge e+v_\ell(q-1)\). At a later support prime it is
already zero because \(q-1\) is an \(\ell\)-adic unit.
At a zero-support prime the integral class-two section has zero
inertia: its two abelian inertia coordinates are zero, and the
remaining pure weight-two inertia is annihilated by the nonzero
integer \(q-1\) in the torsion-free integral coefficient.
The finite reference is unramified there as well. Hence (E.171.3)
reduced modulo \(\ell^{e_\ell}\) is unramified outside
\(\ell\).

For \(\ell\ge5\),
\[
 H^1(G_{\mathbb Q,\{\ell,\infty\}},
                 \mathbb Z/\ell^{e_\ell}(2))=0.
 \tag{E.171.4}
\]
Indeed the integral first group is zero by A.25, and the integral
second group is zero by the weight-two calculation in A.22
(\(K_2(\mathbb Z)\) has no odd torsion). The coefficient exact
sequence gives (E.171.4). Thus the integral section and finite
reference have the same central class modulo \(\ell^{e_\ell}\).

We check the aligning-frame issue, since finite abelian stabilizers
need not vanish. At \(q\ne\ell\), the integral abelian
stabilizer is zero. A finite aligning frame at precision \(N\)
and the reduction of the unique integral abelian aligning frame
differ by an element of
\[
 H^0(\mathbb Q_q,(\mathbb Z/\ell^N)(1)^2)
 \subset \ell^{N-v_\ell(q-1)}
               (\mathbb Z/\ell^N)^2.
 \tag{E.171.5}
\]
The Hall class-two conjugation formulas are integral polynomials.
Changing the abelian part by an element of (E.171.5) changes the
central comparison by a multiple of \(\ell^e\) when
\(N\ge e+v_\ell(q-1)\); its central coordinate may be lifted
arbitrarily. Therefore the finite point match yields zero reduction
of the actual integral local central difference at precision \(e\).
The map
\[
 H^1(\mathbb Q_q,\mathbb Z_\ell(2))
       \longrightarrow H^1(\mathbb Q_q,\mathbb Z/\ell^e(2))
 \tag{E.171.6}
\]
is injective when \(e\ge v_\ell(q^2-1)\): its kernel is
\(\ell^e H^1(\mathbb Q_q,\mathbb Z_\ell(2))\), and the
integral group is \(\mathbb Z_\ell/(q^2-1)\).
This proves the old support matches. At a later support prime
that integral group is zero, so its fixed-abelian class-two fibre
has only one class and is automatically the ordinary point class.

For the dyadic assertion, every odd support prime has
\(v_2(q-1)\le3\) and \(v_2(q^2-1)\le4\): the two private
primes have valuations \((1,4)\) and \((3,4)\), and every
other odd support prime has \((2,3)\). Applying the inertia
calculation to the reference modulo \(2^8\) shows that its
difference from the integral lift is unramified outside two
modulo \(2^5\). E.155 gives
\[
 H^2(G_{\mathbb Q,\{2,\infty\}},\mathbb Z_2(2))\simeq\mathbb Z/2.
 \tag{E.171.7}
\]
In the coefficient exact sequences, reduction from modulus
\(2^5\) to \(2^4\) multiplies the second-cohomology Bockstein
by two. It therefore kills it. The difference modulo sixteen
comes from integral first cohomology, whose global group is
\(\mathbb Z/8\) by A.25 or E.156. One global torsion twist
aligns the reference modulo sixteen. The frame bound
\(8-3\ge4\) and (E.171.6) then prove every required dyadic
support match simultaneously.

At coefficient three all support primes other than three have
\(v_3(q-1)=0\) and \(v_3(q^2-1)=1\).
The difference modulo nine is unramified outside three.
The integral second group at \(\{3,\infty\}\) is zero by
E.66, so the finite difference is the reduction of the integral
\(\mathbb Z/3\) first-cohomology group. One global twist aligns
it, and the same local argument proves all the three-primary
matches. The dyadic and three-primary twists combine into one
\(\mathbb Z/24\) torsor element; there is no independent choice
at different support places. Coefficients dividing no \(q^2-1\)
have zero local central fibre at every off-diagonal support place.

Finally B.25 at class two puts the next obstruction in the
fixed-ramification weight-three group. That group vanishes in
every coefficient: the odd-primary statement is the weight-three
case of A.22, and the dyadic real group is zero by E.155.
Equivalently this is the \(K_4(\mathbb Z)=0\) fixed-ramification
calculation already used in E.137. Thus a global class-three lift
exists, with no assertion that it retains the point matches.
\end{proof}

\begin{result}{E.172}{Proved}{An ordinary arithmetic advance avoiding every previously treated residue}
Fix the point \(t_*=23/64\), its actual class-two quotient with
Hall coordinates modulo \(2^8\), and its actual class-two quotient
modulo nine. Let \(\mathcal P\) be a finite set of primes at
least five. Suppose
\[
 a=23U^8,\quad b=41V^8,\quad c=64W^8
 \tag{E.172.1}
\]
are positive and pairwise coprime, with disjoint new supports,
all of whose primes are five modulo eight and two modulo nine.
The private primes and three are excluded from \(UVW\).
Let \(M\) be divisible by \(43!\) and every prime in
\(\mathcal P\). Assume the exponent-\(M\) diagonal curve
has unit-coordinate points over every finite completion, and
\[
 v_p(a+b-c)\ge2v_p(M)+R
 \quad(p\le B(M)),\qquad B(m)=(m+1)^4,\quad R\ge10.
 \tag{E.172.2}
\]
Suppose the two fixed finite sections have the global abelian
reductions of \((a/c,b/c)\) and match its support point classes
away from their coefficient primes. These classes are independent
of the unit-coordinate point chosen on the curve: the integer
valuations are fixed, and \(M\) kills the relevant two- and
three-primary local unit groups.

For any multiple \(h\) of \(M\), any \(R'\ge10\), and any
additional prime \(\lambda\ge5\), there is an advance
\[
 a'=aP_a^M,\qquad b'=bP_b^M,\qquad c'=cP_c^M
 \tag{E.172.3}
\]
using between eighteen and thirty distinct new primes. They may
avoid any finite set and may all satisfy
\[
 q\equiv5\pmod8,\qquad q\equiv2\pmod9,
 \qquad q\not\equiv\pm1\pmod r
       \quad(r\in\mathcal P\cup\{\lambda\}).
 \tag{E.172.4}
\]
The new exponent-\(h\) curve has unit-coordinate points everywhere,
and (E.172.2) holds for \((a',b',c',h,R')\).
The same two finite global sections still match every support
place, including the newly introduced ones. All retained finite
abelian reductions whose moduli divide \(M\) are unchanged.
\end{result}
\begin{proof}
The finite image input for this reference point is
\[
 z_{t_*}(\ker\chi_{\rm all})=\Phi(Q_2)\times Q_3
 \tag{E.172.5}
\]
in the two specified geometric quotients. This is the proof of
E.127 with inertia at twenty-three and forty-one. Their two-primary
inertia powers act trivially on the two-primary coefficient group;
their section images generate \(Q_2\) by the Frattini criterion.
Commutators of this subgroup give \(Q_2'\) in the cyclotomic
kernel. The two-primary powers of
\([c_\infty,\tau_{23}]\), \([c_\infty,\tau_{41}]\) give
\(-2e_1,-2e_2\) in its abelianization. Together these generate
\(\Phi(Q_2)\). The three-primary commutators give a full
Frattini basis of \(Q_3\), since \(-2\) is a unit.
Primary powers remove the other geometric component.
Conversely the quadratic Kummer characters vanish on the
cyclotomic kernel. This proves (E.172.5). These corrections
remain available over every finite cyclotomic field.

Above a cyclotomic element fixing \(\sqrt{23},\sqrt{41}\),
five modulo eight and two modulo nine, the image therefore
contains the actual unit-one cusp value in each direction.
The infinity cusp's extra quadratic character of \(-1\) is
zero because the element is one modulo four. Thus the required
two-primary difference lies in \(\Phi(Q_2)\); there is no
three-primary restriction. Retain all finite cusp-path marks,
including \(\eta\) and the inverse-root mark of \(-1\).

Choose \(P\ge B(h)\) larger than all old support primes.
At every \(p\le P\), take a unit solution
\((X_p,Y_p,Z_p)\) for exponent \(M\).
At two, five, twenty-three and forty-one choose all coordinates
near one, using (E.172.2) and Newton in a variable with unit
coefficient. The error bound gives a root one modulo eight at
two and one modulo the odd prime at the other places.
The same choice works at every treated prime: it divides
\(M\), hence is at most \(B(M)\).
At three change signs, harmless since \(M\) is even, so every
coordinate is one modulo three.
For each direction choose \(k_i\in\{6,8,10\}\) such that
\[
 2^{k_a}=X_3,\quad2^{k_b}=Y_3,\quad2^{k_c}=Z_3\pmod9.
 \tag{E.172.6}
\]
These even counts give the three residues \(1,4,7\).

Factor each local product target into exactly \(k_i\) units.
Each individual factor must be five modulo eight at two, two
modulo nine at three, a square at both private odd primes,
and nonexceptional at every prime in
\(\mathcal P\cup\{\lambda\}\).
At two the even product is one modulo eight, as required.
At five the allowed residues two and three are nonsquares;
inverse pairs \((2,3)\) realize the target one modulo five.
For \(r\ge7\), every nonzero residue \(u\) is a product
of two residues outside \(\{1,-1\}\): choose the first
outside \(\{1,-1,u,-u\}\). Inverse pairs pad to any larger
even count. At the private primes use their square subgroups,
of orders eleven and twenty; they have nonexceptional inverse
pairs, and the targets are one. This also covers the case that
the active prime is private. At three take the first
\(k_i-1\) factors equal to two and the last equal to the
remaining quotient. At unrestricted places factor arbitrarily.
A principal-unit change of one factor makes each product exactly
the chosen local target. No extraction of a root of order
\(k_i\) in a residue-field unit group is assumed.

Take sufficiently high CRT precisions to make the eventual new
sum error at least \(2v_p(h)+R'\) at every \(p\le P\).
Select the primes successively. For a new prime in direction
\(a,b,c\), with current coefficients \((a_*,b_*,c_*)\),
require respectively that
\[
               c_*/b_*,\qquad c_*/a_*,\qquad -b_*/a_*
 \tag{E.172.7}
\]
be an \(h\)-th power modulo it. Prescribe the actual finite
cusp Frobenius in that direction for the two fixed sections.
Each later prime is also one modulo every earlier new prime.

Let \(F\) contain all the cyclotomic, CRT and finite arithmetic
marks. Only class-two data are retained here: the actual
tangential action on this quotient is determined by \(\chi\),
and its cusp-path values are given by the cyclotomic \(\eta\)
and the indicated roots of \(-1\). Thus their finite arithmetic
fields are contained in a cyclotomic field. No corresponding
assertion about arbitrary higher actual action fields is used.
The prescribed automorphism fixes both private square
roots: reciprocity at twenty-three uses its value one modulo
four, and at forty-one it is the square-residue condition.
It therefore fixes the square root of every ratio \(w\) in
(E.172.7). E.154's private-valuation argument proves
\([F(w^{1/h}):F]=h/2\).
The two fixed section fields over \(F\) form a subgroup of a
product of a two-group and a three-group. Thus their intersection
with the cyclic radical extension can involve only those two
primary parts. The three-part is compatible by E.128 because
\(\chi_3-1\) is a unit. For the two-part the actual cusp value
fixes the marked complementary eighth root; at infinity use
\(\kappa(-b/a)=\kappa(b/c)-\kappa(a/c)-\kappa(-1)\).
The extended E.170 therefore gives the gluing. Its permitted
two-primary conjugation fixes the three-part. All other radical
primary parts have trivial intersection and glue directly.
Compatible primary roots combine to an \(h\)-th root.
Chebotarev now supplies arbitrarily large primes with all the
stated conjugacy and CRT data. The later-prime congruences are
independent CRT conditions; the private valuations and squareclasses
are unchanged as coefficients are updated, so the same argument
applies at every selection.

At each old prime at most \(P\), the product targets solve the
old curve exactly. The imposed precision gives the new sum-error
bound, and Newton at \((1,1,1)\) supplies a unit solution for
exponent \(h\). At a new prime, (E.172.7) supplies the necessary
residue root, which all later prime factors preserve. Simple-root
Hensel in a unit-coefficient variable gives its unit solution;
the derivative is a unit since the new prime exceeds \(B(h)\).
At every remaining prime the coefficients are units and the prime
exceeds \(B(h)\); the diagonal-curve Weil bound of E.31 applies.
This proves the arithmetic invariants.

New exponents are divisible by \(2^9\) and nine, preserving the
fixed finite abelianizations and the old support unit classes.
At a new prime the prescribed reference Frobenius is the actual
unit-one cusp value. Its marked abelian coordinates say that the
current unit parts are the corresponding cusp unit parts.
For coefficient two the local unit group has two-primary exponent
four, since \(q\equiv5\pmod8\); later factors, being
\(2^8\)-th powers, cannot change it. For coefficient three all
local units are three-divisible. The new inertia powers vanish
in the two finite quotients: the dyadic quotient has exponent
dividing \(2^9\), including the infinity peripheral.
Consequently the actual ordinary local point has precisely the
prescribed finite section class. A.31 identifies all subsequent
unit refinements with this same class. This proves the section
assertions and completes the advance.
\end{proof}

\begin{result}{E.173}{Proved}{A coefficient can be installed once without disturbing the ordinary prefix}
Retain an arithmetic prefix of E.172, and let
\(\ell\ge5\) be a prime not yet in \(\mathcal P\).
Let \(\Sigma_0\) be its finite support. Set \(e,u\) as in
(E.171.1) for \(\Sigma_0\), excluding \(q=\ell\), and choose
\[
             N\ge\max(1,e+u),\qquad v_\ell(M)\ge N.
 \tag{E.173.1}
\]
There is an \(A\)-only update
\[
                 a'=a\prod_{q\in T}q^{m_q},\qquad b'=b,\quad c'=c,
 \tag{E.173.2}
\]
with finite new support \(T\), positive integer exponents,
and a finite class-two section \(s_\ell\) modulo \(\ell^N\)
with the following properties. It has abelianization
\((a'/c,b/c)\), is unramified outside
\(\{\ell,\infty\}\cup\Sigma_0\cup T\), and matches the
ordinary support points away from \(\ell\).
Every new prime satisfies (E.172.4) with \(\lambda=\ell\).
The exponent-\(M\) unit curves, the sum-error bound, and both
fixed small-coefficient sections remain valid.
Every \(m_q\) may be required to be divisible by any prescribed
integer prime to \(\ell\), in particular the prime-to-\(\ell\)
part of \(M\) and all previously retained coefficient moduli.

No retained \(\ell\)-section is being refined: this operation
is used only once for a previously untreated coefficient.
Its new exponents need not be globally divisible by \(\ell^N\).
The diagonal-curve invariant is preserved by local power
conditions at old places and a complementary root at new places.
\end{result}
\begin{proof}
Put \(A_0=a/c\), \(B_0=b/c\), and
\(M'=M/\ell^{v_\ell(M)}\). Choose a finite set \(S\)
containing the old support, \(\ell,\infty\), and every prime
at most \(B(M)\). For each finite \(p\in S\), choose a
precision \(d_p\), at least \(2v_p(M)+R\) when
\(p\le B(M)\). Choose an integer \(L\ge N\), as large as
necessary, so that
\[
 L\ge v_\ell(M),\qquad
 L\ge v_\ell\bigl(|(\mathbb Z/p^{d_p})^\times|\bigr)
                 \quad(p\in S, p\text{ finite}).
 \tag{E.173.3}
\]
All these bounds are finite and are fixed before selecting new
primes. Let \(J\), prime to \(\ell\), be divisible by
\(M'\), the additional requested prime-to-\(\ell\) integer,
and the prime-to-\(\ell\) part of every unit-group order in
(E.173.3).

We first construct a permitted auxiliary Frobenius class.
Take a cyclotomic field \(F\) containing \(\mu_{\ell^L}\),
\(\mu_{M'}\), the fixed two-/three-primary arithmetic marks,
and all congruence moduli needed for (E.172.4).
Choose its value five modulo eight, two modulo nine, a square
at twenty-three and forty-one, and nonexceptional at
\(\mathcal P\cup\{\ell\}\). This is possible by CRT.
At the private coefficients the square subgroups still contain
nonexceptional elements; the base has no coefficient-five
quadratic restriction, as proved in E.154.
The finite image argument (E.172.5) supplies the simultaneous
unit-one cusp-zero value for the two reference sections.
Glue the condition that \(w=c/b\) be an \(M'\)-th power,
using precisely the mixed-primary proof of E.172 and E.170.
Let \(D\) be the compositum of these section fields, \(F\),
and the splitting field of \(X^{M'}-w\).
It has an actual automorphism \(\sigma\) with all these
properties. The group \(\operatorname{Gal}(D/F)\) has order
prime to \(\ell\), since the section groups are
two-/three-primary and \(M'\) is prime to \(\ell\).
Since \(\operatorname{Gal}(F/\mathbb Q)\) is abelian,
the commutator subgroup of \(\operatorname{Gal}(D/\mathbb Q)\)
lies in \(\operatorname{Gal}(D/F)\), and therefore has
order prime to \(\ell\).
Enlarge \(S\) by its ramification if needed; only the old
support, primes dividing \(M'\), and \(\ell\) occur, so this
does not change the bounds already described.

Use E.153 for
\(b=\kappa_{\ell^L}(B_0)\) and \(a_0=\kappa_{\ell^L}(A_0)\).
The Kummer class of \(B_0\) is nonzero modulo \(\ell\),
because its private valuation at forty-one is one.
E.154 gives the required independence over \(D\).
At each old support \(q\ne\ell\), prescribe the actual
ordinary point class of a unit solution of the exponent-\(M\)
curve. It has the same full pro-\(\ell\) abelian class as
\((A_0,B_0)\): its unit-coordinate correction is an
\(M\)-th power, and \(v_\ell(M)\ge u\) kills the entire
\(\ell\)-primary local unit group. Its valuations are the
specified old integers.
At a zero-support prime prescribe an unramified \(V_b\)-class
with first coordinate \(a_0\); it exists because both abelian
classes are unramified and the finite-field cohomological
dimension is one. At \(\ell\), a class with that first
coordinate also exists:
\[
 H^2(\mathbb Q_\ell,\mathbb Z/\ell^L(2))=0,
 \tag{E.173.4}
\]
by local duality and the nontrivial Teichm\"uller character on
\(\mathbb Z/\ell^L(-1)\). No point condition at the coefficient
prime is imposed. Thus all the first-coordinate targets are
localizations of the same global \(a_0\), exactly the relative
hypothesis of E.153.

That theorem produces a finite set \(T\) of arbitrarily large
primes with Frobenius class \([\sigma]\) in \(D\), and a
global \(V_b\)-class with all the prescribed old restrictions.
Let \(\bar m_q\in\mathbb Z/\ell^L\) be its new first-coordinate
inertia values. Select positive integer lifts satisfying
\[
                  m_q=\bar m_q\pmod{\ell^L},\qquad J\mid m_q.
 \tag{E.173.5}
\]
CRT permits these choices; when the first residue is zero take
a positive common multiple. The Kummer first coordinate of
(E.173.2) is exactly that of the global finite class. Indeed all
old valuations agree modulo \(\ell^L\), the new valuations
are (E.173.5), and \(-1\) is an \(\ell^L\)-th power over
\(\mathbb Q\). Choose the global geometric frame to use these
literal abelian cocycles and retain the reduction modulo
\(\ell^N\). At the zero-support places the prescribed class
was unramified. Hence the resulting finite section is ramified
only at the asserted old and new support places and at \(\ell\).
At new places E.153 gives the full ordinary point match, not
only its inertia.

We now check the arithmetic invariant, which does not follow
from global \(M\)-divisibility. Write
\(R_0=\prod_{q\in T}q^{m_q}\). At every old finite
\(p\in S\), the local first-coordinate equality gives
\[
                       R_0\in(\mathbb Q_p^\times)^{\ell^L}.
 \tag{E.173.6}
\]
It is a unit. Equation (E.173.5) also makes it a global
\(J\)-th power. In the finite local unit quotient modulo
\(p^{d_p}\), the latter condition kills the prime-to-\(\ell\)
part and (E.173.3)--(E.173.6) kill the \(\ell\)-part.
Consequently \(R_0=1\pmod{p^{d_p}}\). This includes
\(p=\ell\): the factor \(\ell-1\mid J\) removes the
Teichm\"uller part, and an \(\ell^L\)-th power principal unit
is one modulo \(\ell^{L+1}\).
Thus the sum-error bound is preserved.
Moreover (E.173.6), \(M'\mid J\), and the coprimality of
\(M'\) and \(\ell^{v_\ell(M)}\) show that \(R_0\) is a
local \(M\)-th power. Absorb its inverse root into the first
unit coordinate of the old diagonal curve. This preserves the
unit-coordinate point at every old place.

At a new \(q\in T\), the ratio \(c/b\) is an
\(M'\)-th power modulo \(q\). Since \(q-1\) is prime to
\(\ell\), it is also an \(M\)-th power there.
The new first coefficient has positive valuation, and Hensel
therefore gives a unit-coordinate point for exponent \(M\).
The other new primes affect only that first coefficient, so
no cross-prime complementary-root condition arises in this
\(A\)-only operation. All remaining primes exceed \(B(M)\)
and have unit coefficients, so the diagonal Weil bound applies.

Finally \(J\) contains the fixed small-coefficient moduli and
all earlier retained coefficient moduli. The old finite classes
are therefore unchanged globally. Their old local unit classes
are unchanged as well. At a new prime the prescribed reference
Frobenius is geometrically conjugate to cusp zero. Before the
allowed geometric conjugation it fixes both current marked
abelian roots modulo \(2^8\), because the current prefix has
the reference abelianization. After conjugation the two local
Kummer classes are still zero: each Frobenius coordinate changes
by a multiple of \(q-1\). Other auxiliary-prime factors are
\(2^8\)-th powers. Since the local two-primary unit group has
exponent four, its first unit part is indeed that of the
unit-one point; at three all units are divisible. The new
peripheral powers vanish in the retained quotients. This proves
the new small-coefficient point matches and all remaining claims.
\end{proof}

\begin{result}{E.174}{Proved}{An infinite-support ordinary pair with all-prime class-two support matching}
There exists an ordinary completed pair
\[
 (A,B)\in\widehat{\mathbb Q^\times}^{\,2}
 \tag{E.174.1}
\]
with infinite support and one integral all-prime class-two section
whose restriction, at every support prime \(q\), is the actual
ordinary local point class in every coefficient \(\ell\ne q\).
In particular this pair has a global integral class-three lift.
The support can be required to contain \(2,23,41\), with their
valuation pairs respectively
\[
                  (-6,-6),\qquad(1,0),\qquad(0,1),
 \tag{E.174.2}
\]
and every other support prime is five modulo eight and two modulo
nine. At every support prime the two valuations are fixed ordinary
integers, never subsequently changed. The pair is not the image
of a rational point, because its support is infinite.

This theorem supplies no point matching at coefficient-prime
places or at zero-support places, and supplies no class-three
point matching or full pronilpotent lift. It establishes
nonemptiness of the earlier all-coefficient support-matched
class-two gate E.137.
\end{result}
\begin{proof}
Enumerate the primes at least five as
\(\ell_1,\ell_2,\ldots\), with each occurring exactly once.
We construct arithmetic prefixes \((a_n,b_n,c_n,M_n)\), a
finite set of treated coefficients, and one finite reference
section for each treated coefficient. Start with
\[
        (a_0,b_0,c_0)=(23,41,64),\qquad M_0=43!.
 \tag{E.174.3}
\]
The unit-coordinate solution \((1,1,1)\) is valid everywhere,
the sum error is zero, and the two fixed finite sections are
the actual point section at \(23/64\). Thus all initial
hypotheses of E.172 hold. The initial treated set is empty.

Suppose a finite prefix is constructed and the next coefficient
is \(\ell=\ell_n\). Compute \(e,u\) from its present
support by (E.171.1), and choose
\[
                N_\ell\ge\max(1,e+u).
 \tag{E.174.4}
\]
Choose a new exponent \(h\) divisible by the current exponent,
by \((n+43)!\), by \(\ell^{N_\ell}\), and by all
previously retained coefficient moduli. Apply E.172 with this
\(h\), with error parameter at least \(n+10\), and with
\(\ell\) already included among its nonexceptional congruences.
All newly added support primes satisfy
\(q\not\equiv\pm1\pmod\ell\). They therefore contribute
zero to both maxima \(e,u\). This is why (E.174.4) can be
chosen before the advance and does not depend circularly on
its new primes. Put the current exponent equal to \(h\).

Now apply E.173 once, for this coefficient \(\ell\), at
precision \(N_\ell\). Require its correction exponents to
be divisible by the prime-to-\(\ell\) part of the current
exponent and by every earlier retained modulus. Store its
finite section \(s_\ell\), together with its finite support
\(\Sigma_\ell\). Its new support primes are again
nonexceptional at \(\ell\), so (E.174.4) still holds for
the enlarged support. The arithmetic exponent and its local
unit-coordinate invariant remain valid, even though this single
repair need not be globally divisible by \(\ell^{N_\ell}\).
Add \(\ell\) to the treated set. All subsequent advances
and all repairs at other coefficients have exponents divisible
by \(\ell^{N_\ell}\), and all subsequently introduced
primes are nonexceptional at \(\ell\). Every chosen prime
is new, avoids \(2,3,23,41\), and may be larger than any
previously prescribed finite bound.

We verify convergence in the full completion, not merely one
coefficient at a time. Fix a positive integer \(d\). After
finitely many stages, the current factorial exponent is divisible
by \(d\), and every prime divisor of \(d\) which is at
least five has already had its single repair. Every later ordinary
advance therefore has exponents divisible by \(d\). Every
later repair is at a prime not dividing \(d\), and its
prime-to-active-coefficient divisibility likewise includes
\(d\). At two and three there are no repair exceptions.
Consequently
\[
 \text{for every }d,\quad
 d\mid m_q\text{ for all but finitely many newly assigned exponents.}
 \tag{E.174.5}
\]
Each stage uses only finitely many primes. Thus the products of
the assigned prime powers converge in
\(\widehat{\mathbb Q^\times}\). Denote the resulting
coefficient limits by \(a_\infty,b_\infty,c_\infty\),
and put \(A=a_\infty/c_\infty\),
\(B=b_\infty/c_\infty\).

For an ordinary finite prime \(p\), its valuation in each
coefficient is eventually constant: a prime is selected only
once and all later factors are distinct primes. After removing
these fixed powers of \(p\), the coefficient sequences are
Cauchy in \(\mathbb Z_p^\times\). Indeed, for any fixed
precision \(p^k\), taking \(d\) divisible by the exponent
of \((\mathbb Z/p^k)^\times\) in (E.174.5) makes every
later unit multiplier one modulo \(p^k\). Their limits are
nonzero. Every fixed \(p\) eventually lies below \(B(M_n)\),
and the increasing sum-error bound is preserved by every repair.
Passing to these actual \(p\)-adic limits gives
\[
                 a_{\infty,p}+b_{\infty,p}=c_{\infty,p}.
 \tag{E.174.6}
\]
Thus \(t_p=a_{\infty,p}/c_{\infty,p}\) and
\(1-t_p=b_{\infty,p}/c_{\infty,p}\) are both nonzero,
and give the required actual ordinary local point. At the real
place all coefficient signs are positive, so the point \(1/2\)
has the required completed multiplicative images. This proves
ordinarity at every place. Each ordinary advance adds at least
eighteen fresh support primes. Hence support is infinite, while
all its valuations remain the positive integers assigned in one
of the three directions, or the negative common integer in the
denominator direction. The initial private valuations are exactly
(E.174.2).

We next verify the finite references against these final points.
For coefficients two and three, every new exponent is divisible
by \(2^9\) and nine. The final global abelian reductions are
therefore those of \(23/64\). At every old support place the
full two-/three-primary local unit classes are unchanged; at a
new support place E.172 or E.173 gave the actual finite cusp
match, including its unit part and Frobenius. Subsequent factors
are powers of the retained moduli. Since the relevant local
unit groups have two-primary exponent at most eight at the
private primes and four at fresh primes, and have trivial
three-primary part, these changes do not alter their full local
unit classes. The fixed valuations and A.31 therefore identify
the final actual local point with the finite point class used in
that construction. Both fixed finite reference sections satisfy
the small-coefficient hypotheses of E.171 at every support place.

Fix now \(\ell\ge5\), and consider the stage at which its
section \(s_\ell\) was stored. At every old support place
\(q\ne\ell\), the ordinary point target came from an
exponent-\(M\) unit-coordinate solution with
\(v_\ell(M)\ge N_\ell\ge u\). Its unit correction is
therefore trivial in the full pro-\(\ell\) local Kummer
group. During the repair, the first-coordinate relative condition
keeps the old full local unit class: its quotient is an
\(\ell^L\)-th power with \(L\ge N_\ell\ge u\),
and its valuations do not change. At the repair's new support
places \(q-1\) is prime to \(\ell\), so every unit class
is already trivial. All subsequent updates have exponents
divisible by \(\ell^{N_\ell}\), which continue to kill
the old local unit groups. Hence the final ordinary point has
the same full pro-\(\ell\) abelian boundary data as the
stored point target at every \(q\in\Sigma_\ell\),
\(q\ne\ell\). By A.31 its point-section class there is
the same in every nilpotent quotient, in particular in the finite
class-two quotient under discussion.

The global abelian reduction of the final pair is the one of
\(s_\ell\), because all later exponents are divisible by
\(\ell^{N_\ell}\). This reference is unramified outside
\(\{\ell,\infty\}\cup\Sigma_\ell\). Every support
prime added afterwards is different from \(\ell\) and is
nonexceptional at \(\ell\). Thus every hypothesis of the
first part of E.171 is fulfilled. It proves that the integral
class-two section matches the final actual point at every
off-diagonal support place for this coefficient. This argument
applies to every \(\ell\ge5\); each is treated once and
its finite verification remains valid forever.

Finally E.171 supplies the one simultaneous global
\(\mathbb Z/24\) twist for the two small coefficients and
assembles all prime coefficients into one integral class-two
section. The same theorem, using B.25 and the fixed-ramification
weight-three vanishing, gives its global integral class-three
lift. These arguments use no coefficient-prime point condition,
no matching at zero-support places, and no higher-layer matching
assumption. The stated limits on the conclusion follow.

The finite certificate \path{code/ghost_audit_all_prime_class_two.py},
with saved output \path{code/ghost_audit_all_prime_class_two.out},
checks the nonexceptional residue factorizations, private quadratic
and even-count shells, simultaneous CRT data, finite inertia/frame
precision, dyadic Hall exponent, and the relative unit-power
intersection used above. It does not compute the Chebotarev or
Poitou--Tate choices and is not a finite substitute for this
infinite construction.
\end{proof}
